\ifdefined\gkver\else\def\gkver{student}\fi
\documentclass[\gkver]{gaokaozhenti}
\begin{document}

\chapter{2005年全国理综Ⅱ}
%% paperType: 理综物理部分

\begin{enumerate}
\item
%% number: 14
%% paperName: 2005年全国理综Ⅱ第 14 题 3 分
%% typeId: 1
%% type: 单选题
%% chapter: 运动和力
%% point: 牛顿第二定律
%% method: 无
%% score: 3
%% degree: 650
%% duplicateId: 0
%% body:
如图所示，位于光滑固定斜面上的小物块 P 受到一水平向右的推力 $F$ 的作用。已知物块 P 沿斜面加速下滑。现保持 $F$ 的方向不变，使其减小，则加速度\xzanswer{B}

\onepicture[w=4.5cm]{figs/14.png}

\fourchoices[answer=B]
{一定变小}
{一定变大}
{一定不变}
{可能变小，可能变大，也可能不变}

%% answer: B
%% memo:
\memoanswer{
本题原卷及材料中未提供详解，待补充。
}

\item
%% number: 15
%% paperName: 2005年全国理综Ⅱ第 15 题 3 分
%% typeId: 2
%% type: 多选题
%% chapter: 光学
%% point: 光的折射
%% method: 无
%% score: 3
%% degree: 650
%% duplicateId: 0
%% body:
一束复色光由空气射向玻璃，发生折射而分为 a、b 两束单色光，其传播方向如图所示。设玻璃对 a、b 的折射率分别为 $n_{a}$ 和 $n_{b}$，a、b 在玻璃中的传播速度分别为 $v_{a}$ 和 $v_{b}$，则\xzanswer{AD}

\onepicture[w=4.5cm]{figs/15.png}

\fourchoices[answer=AD]
{$n_{a}>n_{b}$}
{$n_{a}<n_{b}$}
{$v_{a}>v_{b}$}
{$v_{a}<v_{b}$}

%% answer: AD
%% memo:
\memoanswer{
本题原卷及材料中未提供详解，待补充。
}

\item
%% number: 16
%% paperName: 2005年全国理综Ⅱ第 16 题 3 分
%% typeId: 2
%% type: 多选题
%% chapter: 气体性质
%% point: 热力学第一定律
%% method: 无
%% score: 3
%% degree: 650
%% duplicateId: 0
%% body:
对于定量气体，可能发生的过程是\xzanswer{AC}

\fourchoices[answer=AC]
{等压压缩，温度降低}
{等温吸热，体积不变}
{放出热量，内能增加}
{绝热压缩，内能不变}

%% answer: AC
%% memo:
\memoanswer{
本题原卷及材料中未提供详解，待补充。
}

\item
%% number: 17
%% paperName: 2005年全国理综Ⅱ第 17 题 3 分
%% typeId: 1
%% type: 单选题
%% chapter: 光学
%% point: 光电效应
%% method: 无
%% score: 3
%% degree: 650
%% duplicateId: 0
%% body:
图中画出了氢原子的 4 个能级，并注明了相应的能量 $E$。处在 $n=4$ 的能级的一群氢原子向低能级跃迁时，能够发出若干种不同频率的光波。已知金属钾的逸出功为 $2.22\UeV$。在这些光波中，能够从金属钾的表面打出光电子的总共有\xzanswer{C}

\onepicture[h=4.5cm]{figs/17.png}

\fourchoices[answer=C]
{二种}
{三种}
{四种}
{五种}

%% answer: C
%% memo:
\memoanswer{
本题原卷及材料中未提供详解，待补充。
}

\item
%% number: 18
%% paperName: 2005年全国理综Ⅱ第 18 题 3 分
%% typeId: 2
%% type: 多选题
%% chapter: 曲线运动
%% point: 天体运动
%% method: 近似与估算
%% score: 3
%% degree: 450
%% duplicateId: 0
%% body:
已知引力常量 $G$、月球中心到地球中心的距离 $R$ 和月球绕地球运行的周期 $T$。仅利用这三个数据，可以估算出的物理量有\xzanswer{BD}

\fourchoices[answer=BD]
{月球的质量}
{地球的质量}
{地球的半径}
{月球绕地球运行速度的大小}

%% answer: BD
%% memo:
\memoanswer{
本题原卷及材料中未提供详解，待补充。
}

\item
%% number: 19
%% paperName: 2005年全国理综Ⅱ第 19 题 3 分
%% typeId: 2
%% type: 多选题
%% chapter: 机械波
%% point: 波的图象
%% method: 无
%% score: 3
%% degree: 450
%% duplicateId: 0
%% body:
一简谐横波沿 $x$ 轴正方向传播，某时刻其波形如图所示。下列说法正确的是\xzanswer{AD}

\onepicture[w=6.5cm]{figs/19.png}

\fourchoices[answer=AD]
{由波形图可知该波的波长}
{由波形图可知该波的周期}
{经 1/4 周期后质元 P 运动到 Q 点}
{经 1/4 周期后质元 R 的速度变为零}

%% answer: AD
%% memo:
\memoanswer{
本题原卷及材料中未提供详解，待补充。
}

\item
%% number: 20
%% paperName: 2005年全国理综Ⅱ第 20 题 3 分
%% typeId: 1
%% type: 单选题
%% chapter: 交流电
%% point: 交流电
%% method: 无
%% score: 3
%% degree: 650
%% duplicateId: 0
%% body:
处在匀强磁场中的矩形线圈 abcd，以恒定的角速度绕 ab 边转动，磁场方向平行于纸面并与 ab 垂直。在 $t=0$ 时刻，线圈平面与纸面重合（如图），线圈的 cd 边离开纸面向外运动。若规定由 a$\to$b$\to$c$\to$d$\to$a 方向的感应电流为正，则能反映线圈中感应电流 $I$ 随时间 $t$ 变化的图线是\xzanswer{C}

\onepicture[w=4.5cm]{figs/20a.png}

\fourchoices[answer=C, ispicture=true, w=3cm]
{figs/20b.png}
{figs/20c.png}
{figs/20d.png}
{figs/20e.png}

%% answer: C
%% memo:
\memoanswer{
本题原卷及材料中未提供详解，待补充。
}

\item
%% number: 21
%% paperName: 2005年全国理综Ⅱ第 21 题 3 分
%% typeId: 2
%% type: 多选题
%% chapter: 电场
%% point: 电场叠加
%% method: 无
%% score: 3
%% degree: 450
%% duplicateId: 0
%% body:
图中 a、b 是两个点电荷，它们的电量分别为 $Q_{1}$、$Q_{2}$，MN 是 ab 连线的中垂线，P 是中垂线上的一点。下列哪种情况能使 P 点场强方向指向 MN 的左侧？\xzanswer{ACD}

\onepicture[w=4.5cm]{figs/21.png}

\fourchoices[answer=ACD]
{$Q_{1}$、$Q_{2}$ 都是正电荷，且 $Q_{1}<Q_{2}$}
{$Q_{1}$ 是正电荷，$Q_{2}$ 是负电荷，且 $Q_{1}>|Q_{2}|$}
{$Q_{1}$ 是负电荷，$Q_{2}$ 是正电荷，且 $|Q_{1}|<Q_{2}$}
{$Q_{1}$、$Q_{2}$ 都是负电荷，且 $|Q_{1}|>|Q_{2}|$}

%% answer: ACD
%% memo:
\memoanswer{
本题原卷及材料中未提供详解，待补充。
}

\item
%% number: 22
%% paperName: 2005年全国理综Ⅱ第 22 题 12 分
%% typeId: 4
%% type: 实验题
%% chapter: 电路
%% point: 测电动势与内阻
%% method: 无
%% score: 12
%% degree: 450
%% duplicateId: 0
%% body:
\begin{enumerate}
	\item
	用游标为 50 分度的卡尺（测量值可准确到 $0.02\Umm$）测定某圆柱的直径时，卡尺上的示数如图。可读出圆柱的直径为 \tkanswer[2cm]{42.12}~$\Umm$。

	\onepicture[w=6cm]{figs/22.png}

	\item
	利用图 1 所示的电路测量电流表 mA 的内阻 $R_{A}$。图中 $R_{1}$、$R_{2}$ 为电阻，$K_{1}$、$K_{2}$ 为电键，B 是电源（内阻可忽略）。

	\begin{steps}
		\item
		根据图 1 所给出的电路原理图，在图 2 的实物图上连线。
		\item
		已知 $R_{1}=140\UO$，$R_{2}=60\UO$。当电键 $K_{1}$ 闭合、$K_{2}$ 断开时，电流表读数为 $6.4\UmA$；当 $K_{1}$、$K_{2}$ 均闭合时，电流表读数为 $8.5\UmA$。由此可以求出 $R_{A}=$ \tkanswer[2cm]{43} $\UO$。（保留 2 位有效数字）
	\end{steps}
\end{enumerate}

\onepicture[w=4.5cm, align=right]{figs/22a.png}

\drawpicanswer{%
\onepicture[w=5.5cm, align=right]{figs/22b.png}%
}{%
\onepicture[w=5.5cm, align=right]{figs/22_an.png}%
}

%% answer:
\jdanswer{
\begin{enumerate}
	\item
	$42.12\Umm$。
	\item
	①实物连线如图所示。

	②$R_{A}=43\UO$。
\end{enumerate}
}

%% memo:
\memoanswer{
本题原卷及材料中未提供详解，待补充。
}

\item
%% number: 23
%% paperName: 2005年全国理综Ⅱ第 23 题 16 分
%% typeId: 6
%% type: 计算题
%% chapter: 功和能
%% point: 动能定理
%% method: 整体与隔离
%% score: 16
%% degree: 650
%% duplicateId: 0
%% body:
如图所示，在水平桌面的边角处有一轻质光滑的定滑轮 K，一条不可伸长的轻绳绕过 K 分别与物块 A、B 相连，A、B 的质量分别为 $m_{A}$、$m_{B}$。开始时系统处于静止状态。现用一水平恒力 $F$ 拉物块 A，使物块 B 上升。已知当 B 上升距离为 $h$ 时，B 的速度为 $v$。求此过程中物块 A 克服摩擦力所做的功。重力加速度为 $g$。

\onepicture[h=4.2cm, align=right]{figs/23.png}

%% answer:
\jdanswer{$W=Fh-m_{B}gh-\frac{1}{2}(m_{A}+m_{B})v^{2}$}

%% memo:
\memoanswer{
由于连结 AB 绳子在运动过程中未松，故 AB 有一样的速度大小，对 AB 系统，由功能关系有：

$Fh-W-m_{B}gh=\frac{1}{2}(m_{A}+m_{B})v^{2}$

求得：$W=Fh-m_{B}gh-\frac{1}{2}(m_{A}+m_{B})v^{2}$。
}

\item
%% number: 24
%% paperName: 2005年全国理综Ⅱ第 24 题 19 分
%% typeId: 6
%% type: 计算题
%% chapter: 磁场
%% point: 带电粒子在磁场中的运动
%% method: 分情况讨论
%% score: 19
%% degree: 550
%% duplicateId: 0
%% body:
在同时存在匀强电场和匀强磁场的空间中取正交坐标系 Oxyz（$z$ 轴正方向竖直向上），如图所示。已知电场方向沿 $z$ 轴正方向，场强大小为 $E$；磁场方向沿 $y$ 轴正方向，磁感应强度的大小为 $B$；重力加速度为 $g$。问：一质量为 $m$、带电量为 $+q$ 的从原点出发的质点能否在坐标轴（$x$，$y$，$z$）上以速度 $v$ 做匀速运动？若能，$m$、$q$、$E$、$B$、$v$ 及 $g$ 应满足怎样的关系？若不能，说明理由。

\onepicture[h=4cm, align=right]{figs/24.png}

%% answer:
\jdanswer{
能。

第一种情况：$mg>qE$ 时，由平衡条件知洛伦兹力 $f$ 沿 $z$ 轴正向，粒子以 $v$ 沿 $x$ 轴正向运动，由匀速运动易知其条件是 $mg-qE=qvB$；

第二种情况：$mg<qE$ 时，洛伦兹力 $f$ 沿 $z$ 轴负方向，粒子以 $v$ 沿 $x$ 轴负向运动，由匀速运动知条件是 $qE-mg=qvB$。
}

%% memo:
\memoanswer{
能。

第一种情况：$mg>qE$，由平衡条件知洛伦兹力 $f$ 沿 $z$ 轴正向，粒子以 $v$ 沿 $x$ 轴正向运动由匀速运动易知其条件是：$mg-qE=qvB$。

第二种情况：$mg<qE$，则 $f$ 沿 $z$ 轴负方向，粒子以 $v$ 沿 $x$ 轴负向运动，由匀速运动知条件是：$qE-mg=qvB$。
}

\item
%% number: 25
%% paperName: 2005年全国理综Ⅱ第 25 题 20 分
%% typeId: 6
%% type: 计算题
%% chapter: 运动和力
%% point: 动量和能量
%% method: 无
%% score: 20
%% degree: 450
%% duplicateId: 0
%% body:
质量为 $M$ 的小物块 A 静止在离地面高 $h$ 的水平桌面的边缘，质量为 $m$ 的小物块 B 沿桌面向 A 运动以速度 $v_{0}$ 与之发生正碰（碰撞时间极短）。碰后 A 离开桌面，其落地点离出发点的水平距离为 $L$。碰后 B 反向运动。求 B 后退的距离。已知 B 与桌面间的动摩擦因数为 $\mu$。重力加速度为 $g$。

\onepicture[h=4cm, align=right]{figs/25.png}

%% answer:
\jdanswer{$s=\frac{1}{2\mu g}\left(\frac{ML}{m}\sqrt{\frac{g}{2h}}-v_{0}\right)^{2}$}

%% memo:
\memoanswer{
设 AB 碰后 A 的速度为 $v_{1}$，则 A 平抛有：$h=\frac{1}{2}gt^{2}$，$L=v_{1}t$

求得：$v_{1}=L\sqrt{\frac{g}{2h}}$ ①

设碰后 B 的速度为 $v_{2}$，则对 AB 碰撞过程由动量守恒有：$mv_{0}=Mv_{1}-mv_{2}$ ②

设 B 后退距离为 $s$，对 B 后退直至停止过程，由动能定理：$\mu mgs=\frac{1}{2}mv_{2}^{2}$ ③

由①②③解得：$s=\frac{1}{2\mu g}\left(\frac{ML}{m}\sqrt{\frac{g}{2h}}-v_{0}\right)^{2}$。
}

\end{enumerate}

\end{document}
